\documentclass[10pt,a4paper]{amsart}
\usepackage[margin=19mm]{geometry}
\usepackage{amsmath,amssymb,mathtools}
\usepackage[scr=rsfs,cal=euler]{mathalfa}
\usepackage{xfrac}
\usepackage[unicode,hidelinks,bookmarks=false]{hyperref}
\newcommand{\ie}{i.e.\ }
\newcommand{\cf}{cf.\ }
\newcommand{\resp}{respectively\ }
\newcommand{\Subsection}{\subsection}
\providecommand{\nobreakdash}{\nobreak\hbox{-}}
\setlength{\emergencystretch}{2em}
\multlinegap0pt
\usepackage{bm}
\usepackage{bbm}
\usepackage{dsfont}
\usepackage{xspace}
\usepackage{cancel}
\usepackage{mathtools}
\usepackage{amsthm}
\makeatletter
\@ifundefined{theorem}{\newtheorem{theorem}{Theorem}}{}
\@ifundefined{assumption}{\newtheorem{assumption}[theorem]{Assumption}}{}
\@ifundefined{lemma}{\newtheorem{lemma}[theorem]{Lemma}}{}
\@ifundefined{proposition}{\newtheorem{proposition}[theorem]{Proposition}}{}
\makeatother
\title[Integrated density of states for the Poisson point interacti]{Integrated density of states for the Poisson point interactions on $\mathbf{R}^3$}
\author{Masahiro Kaminaga, Takuya Mine, Fumihiko Nakano}
\date{Source: arXiv:2406.02256v2 (CC BY 4.0)}
\begin{document}
\maketitle

\noindent\emph{Source and licence.} Integrated density of states for the Poisson point interactions on $\mathbf{R}^3$. Version arXiv:2406.02256v2 (CC BY 4.0), distributed under the Creative Commons Attribution 4.0 International licence, as recorded in the arXiv licence metadata for this exact version and shown on the abstract page https://arxiv.org/abs/2406.02256v2. Full rights notice: licenses/arxiv-2406.02256-CC-BY-4.0.txt.\par\smallskip

\noindent\textit{Editorial scope.} Counted unit: the Lemma whose source label is \texttt{lemma\_ids-est} in the article (source0.tex lines 907--933, source label \texttt{lemma\_ids-est}), delivered together with the article's own complete original proof of it (source0.tex lines 938--1052, one \texttt{proof} environment of 115 lines). No mathematics is written, repaired or re-derived: statement and proof are the article's own text line for line. The proof is detached from the statement in the source: the article states the result at lines 907--933 and prints its proof at lines 938--1052, and the proof environment's own header names the statement, so the association is the article's own. Required context retained uncounted: assumption\_PA (lines 223--238); theorem\_asymptotics\_IDS (lines 467--493); 01\_08 (lines 482--486); prop\_ev2 (lines 717--749); 03\_01 (lines 798--801); prop\_cluster (lines 802--819). Every definition, display and label of those lines is retained unchanged. Omitted from this package and not redistributed: 1--222, 239--466, 487--716, 750--797, 820--906, 934--937, 1053--5445. Nothing inside a retained excerpt is omitted or altered: every retained body is delivered exactly as the article's lines are written, so no omission receipt of any kind accompanies this package. Citations: the retained excerpts carry no citation command at all, so no bibliography entry of the article is transcribed and no reference is redistributed. Reference adaptation: none was needed and none was made; every reference inside the delivered unit resolves inside it, so no unresolved reference or citation is produced. Printed slips kept literal: no printed word, symbol, hypothesis or number of the counted statement or of its proof was altered. Layout and class: the article's own class is \texttt{article} with packages including amsmath, amssymb, amsthm, graphicx, color, pict2e; the wrapper is the lane's pinned \texttt{amsart} editorial wrapper at 10pt on A4, which restates the theorem-like environments the retained text needs and adds the article's own macro definitions that the retained text needs (none), with extra wrapper packages (bm, bbm, dsfont, xspace, cancel, mathtools). The delivered numbering is the unit's own. Assets: no retained span contains a figure environment, a graphics-inclusion command, a PDF-page inclusion, a table or a TikZ picture; no image file is copied into this package, the package declares no asset dependency and no image file is redistributed with it. Licence: version arXiv:2406.02256v2 of this article is distributed under the Creative Commons Attribution 4.0 International licence, as recorded in the arXiv licence metadata for this exact version and shown on the abstract page https://arxiv.org/abs/2406.02256v2. A full rights notice is delivered at licenses/arxiv-2406.02256-CC-BY-4.0.txt. Characters: every retained excerpt is pure ASCII, so the delivered engine, XeLaTeX with its default Latin Modern text font, reports no missing character.
\begin{assumption}
\label{assumption_PA}
\begin{enumerate}
 \item 
The space $(\Omega,\mathcal{F},\mathbf{P})$ is a probability space.
The random set $Y_\omega$ ($\omega\in \Omega$) is the Poisson configuration
(the support of the Poisson point process) on $\mathbf{R}^d$
with intensity measure $\rho dx$ for some constant $\rho>0$.

 \item The parameters $\alpha_\omega=(\alpha_{\omega,y})_{y\in Y_\omega}$
are real-valued i.i.d.\ random variables
with common distribution $\nu$ on $\mathbf{R}$.
Moreover, $(\alpha_{\omega,y})_{y\in Y_\omega}$ 
are independent of $Y_\omega$.
\end{enumerate}
\end{assumption}

\begin{theorem}
 \label{theorem_asymptotics_IDS}
Let $d=3$.
Suppose that $\alpha_\omega$ and $Y_\omega$ satisfy
Assumption \ref{assumption_PA}.
Suppose further that
$\alpha_\omega$ is a constant sequence
with its common value $\alpha$.
%Suppose 
%and $Y_\omega$ is a Poisson configuration with intensity $\rho dx$,
%where $\rho$ is a positive constant.
Let $N(\lambda)$ be the corresponding IDS
for the operator $H_\omega=-\Delta_{\alpha,Y_\omega}$.
Then, for any $\epsilon$ with $0<\epsilon<3$,
we have
\begin{align}
\label{01_08}
N(-s^2)=\frac{2\pi}{3}\rho^2 R_\alpha(s)^3 + O(s^{-6+\epsilon})
\quad(s\to \infty).
\end{align}
In particular, the principal term of the asymptotics is given by
\begin{align}
\label{01_09}
 N(-s^2) \sim \frac{2\pi}{3}\rho^2 t_0^3s^{-3}\quad
(s\to\infty).
\end{align}
\end{theorem}

\begin{proposition}
\label{prop_ev2}
Let 
$t_0$ be the unique positive solution of
$t-e^{-t}=0$. 
\begin{enumerate}
  \item[(1)] Suppose $\alpha=0$.
Then we have
$s_0(R)=t_0/R$
and
$R_0(s)=t_0/s$.



  \item[(2)] Suppose $\alpha\not=0$.
Then, for any $\epsilon>0$,
there exists a positive constant $R_0$ dependent on $\alpha$ and
$\epsilon$
such that
\begin{align}
\label{02_03}
\frac{t_0-\epsilon}{R}<s_\alpha(R)<\frac{{t_0+\epsilon}}{R}
\end{align}
for any $R$ with $0<R<R_0$,
and 
\begin{align}
\label{02_04}
 \frac{t_0-\epsilon}{s}<
R_\alpha(s)<\frac{t_0+\epsilon}{s}
\end{align}
for any $s$ with $s>(t_0+\epsilon)/R_0$.
 \end{enumerate}
\end{proposition}

\begin{align}
\label{03_01}
n_y(R) :=\#(Y_\omega \cap B_y(R)).
\end{align}

\begin{proposition}
\label{prop_cluster}
Let $Y_\omega$ be the Poisson configuration 
with intensity $\rho dx$, where $\rho$ is a positive constant.
Let $L\in\mathbf{N}$ and put $Q_L:=(0,L)^d$.
Then, we have for $n=1,2,3,\ldots$ and $R>0$
\begin{align}
\label{03_02}
 \frac{\mathbf{E}[
\#\{y\in Y_\omega \cap Q_L
\,;\, n_y(R)=n
\}
]}
{|Q_L|}
=
\frac{1}{(n-1)!}|B_0(R)|^{n-1}\rho^n e^{-\rho|B_0(R)|}.
\end{align}
\end{proposition}

\begin{lemma}
\label{lemma_ids-est}
 For every $\delta$ with $1/2<\delta <1$,
and for every $m>0$,
there exist constants $R_0=R_0(\rho,\alpha,\delta,m)>0$
and $C=C(\rho,\alpha,\delta,m)>0$
such that
\begin{align}
& \mathbf{E}[N_{\alpha_\omega,Y_\omega,Q_L}(-(s_\alpha(R)-R^m)^2)]\notag\\
\geq&\  \frac{1}{2}
\mathbf{E}[
\#\{y\in Y_\omega\cap Q_L\,;\, n_y(R)=2\}]
-
C R^{6\delta}|Q_L|,
\label{04_01}
\\
& \mathbf{E}[N_{\alpha_\omega,Y_\omega,Q_L}(-(s_\alpha(R)+R^m)^2)]\notag\\
\leq&\  \frac{1}{2}
\mathbf{E}[
\#\{y\in Y_\omega\cap Q_L\,;\, n_y(R)=2\}]
+
C R^{6\delta}|Q_L|,
\label{04_02}
\end{align}
for every $0<R<R_0$ and every $L>R^{-2}$,
where $n_y(R)$ is defined in (\ref{03_01}).
\end{lemma}

\begin{proof}[Proof of Theorem \ref{theorem_asymptotics_IDS}]
Assume Lemma \ref{lemma_ids-est} holds.
Put $f(R)=s_\alpha(R)-R^m$.
By Proposition \ref{prop_ev2},
we see that $s_\alpha(R)\sim t_0/R$ ($R\to 0$) and
$R_\alpha(s)\sim t_0/s$ ($s\to \infty$).
Thus $f(R)\sim s_\alpha(R)\sim t_0/R$ ($R\to 0$).


Since 
the function $R=R_\alpha(s)$ is the inverse of 
the function $s=s_\alpha(R)$, we have
$R=R_\alpha(f(R)+R^m)$.
By the Taylor theorem, we have
\begin{align}
R=
R_\alpha(f(R))+\frac{d R_\alpha}{d s}(f(R)+\theta R^m)R^m,\quad 0<\theta<1.
\label{04_03}
\end{align}
In order to calculate the derivative $dR_\alpha/ds$,
recall that $R_\alpha$ is defined as the solution of the 
following equation with respect to $R$
\begin{align*}
 h(s,R)=-4\pi\alpha,\quad
\mbox{where }
h(s,R):=s- \frac{e^{-sR}}{R}.
\end{align*}
Then
\begin{align*}
 \frac{\partial h}{\partial s}=
1+e^{-sR}>0,\quad
 \frac{\partial h}{\partial R}=
\frac{(sR+1)e^{-sR}}{R^2}>0,
\end{align*}
and we have by the implicit function theorem
\begin{align}
&\frac{dR_\alpha}{ds}(s)
=-
\left.
\frac{\frac{\partial h}{\partial s}}{\frac{\partial h}{\partial R}}
\right|_{R=R_\alpha(s)}
=
-\frac{ 1+e^{-sR_\alpha(s)}}{(sR_\alpha(s)+1)e^{-sR_\alpha(s)}}R_\alpha(s)^2,
\label{04_04}\\
&\frac{ds_\alpha}{dR}(R)
=-
\left.
\frac{\frac{\partial h}{\partial R}}{\frac{\partial h}{\partial s}}
\right|_{s=s_\alpha(R)}
=
-
\frac{(s_\alpha(R)R+1)e^{-s_\alpha(R)R}}{ 1+e^{-s_\alpha(R)R}}R^{-2}.
\label{04_05}
\end{align}
Then there exist positive constants 
$C_1$, $C_2$, $C_3$, $C_4$ such that
\begin{align*}
&\frac{C_1}{R}\leq f(R)+\theta R^m \leq \frac{C_2}{R},\\
&C_3 R\leq R_\alpha(f(R) + \theta R^m)\leq C_4 R
\end{align*}
for every $0\leq \theta \leq 1$ and small $R>0$.
Thus (\ref{04_03}) and (\ref{04_04}) imply
\begin{align}
\label{04_06}
 R = R_\alpha(f(R)) + O(R^{m+2})
=R_\alpha(f(R))(1+O(R^{m+1}))
\quad (R\to 0).
\end{align}
Take $m\geq 2$.
Then, by (\ref{04_01}), (\ref{04_06}) and Proposition \ref{prop_cluster}
\begin{align*}
 \frac{\mathbf{E}[N_{\alpha_\omega,Y_\omega,Q_L}(-f(R)^2)]}{|Q_L|}
\geq&\ 
\frac{2\pi}{3}R^3\rho^2(1+O(R^3))
- C R^{6\delta}\\
\geq &\ 
\frac{2\pi\rho^2}{3}R_\alpha(f(R))^3 - C' f(R)^{-6\delta}
\end{align*}
for some $C'>0$,
 and sufficiently small $R$ and large $L$.
Moreover,
we have by (\ref{04_05})
\begin{align*}
\frac{df}{dR}(R)
=&\ 
-
\frac{(s_\alpha(R)R+1)e^{-s_\alpha(R)R}}{ 1+e^{-s_\alpha(R)R}}R^{-2}
 - mR^{m-1}\\
\sim&\ 
-\frac{(t_0+1)e^{-t_0}}{ 1+e^{-t_0}}R^{-2}
=-t_0 R^{-2}
<0
\end{align*}
(notice that $t_0=e^{-t_0}$) for sufficiently small $R>0$.
Thus 
there exists some $R_0>0$ such that
the map $s=f(R)$ is 
bijective from $(0,R_0)$ to $(f(R_0),\infty)$,
and the lower bound in (\ref{01_08}) is proved.


Similarly, in order to prove the upper bound in (\ref{01_08}),
we put $g(R)=s_\alpha(R)+R^m$.
Then we have from (\ref{04_02}) and Proposition \ref{prop_cluster}
\begin{align*}
\frac{\mathbf{E}[N_{\alpha_\omega,Y_\omega,Q_L}(-g(R)^2)]}{|Q_L|}
\leq &\ 
\frac{2\pi\rho^2}{3}R_\alpha(g(R))^3 + C'' g(R)^{-6\delta}
\end{align*}
for some $C''>0$,
 and sufficiently small $R$ and large $L$.
Making the change of variable $s=g(R)$,
and taking the limit $L\to \infty$,
we have the conclusion of Theorem \ref{theorem_asymptotics_IDS}.
\end{proof}

\end{document}
